\medskip





\medskip

\begin{enumerate}
\item On considère la fonction $g$ définie sur $\R$ par : 

\[ g(x) = \dfrac{2\text{e}^x}{\text{e}^x +1}.\]

\smallskip









On a pour tout réel $x$, \: $g(x) = \dfrac{\text{e}^x \times 2 }{\text{e}^x\left(1 + \text{e}^{-x}\right)} = \dfrac{2}{1 + \text{e}^{-x}}$.

Or $\displaystyle\lim_{x \to + \infty} \text{e}^{-x} = 0$, donc $\displaystyle\lim_{x \to + \infty} 1 + \text{e}^{-x} = 1$ et enfin $\displaystyle\lim_{x \to + \infty} g(x) = \dfrac{2}{1} = 2$. \hfill Réponse \textbf{b.}

\item La dérivée seconde $f''$ est positive sur les intervalles $]- \infty~;~-1[$ et $[2~;~+ \infty[$, donc la fonction $f$ est convexe sur ces intervalles. \hfill Réponse \textbf{c.}

\item 









On a pour tout naturel $n$ :

$b_{n+1} = a_{n+1} - 2 = 0,5a_n + 1 - 2 = 0,5a_n - 1 = 0,5\left(a_n - 2 \right) = 0,5b_n$.

L'égalité $b_{n+1} = 0,5b_n$ montre que la suite $\left(b_n\right)$ est une suite géométrique de raison 0,5. 

\hfill Réponse \textbf{d.}

\item 











$\bullet~~$Comme $- 1 < \dfrac{1}{4} < 1$, on sait que $\displaystyle\lim_{n \to + \infty}\left(\dfrac14\right)^n = 0$ et donc $\displaystyle\lim_{n \to + \infty}1 + \left(\dfrac14\right)^n = 1$ ;

$\bullet~~$On a $\dfrac{n}{n + 1} = \dfrac{1}{1 + \frac1 n}$.
Or $\displaystyle\lim_{n \to + \infty} \frac1 n = 0 \Rightarrow \displaystyle\lim_{n \to + \infty} \dfrac{n}{n + 1} = 1$ et $\displaystyle\lim_{n \to + \infty}2 - \dfrac{n}{n + 1} = 1$.



On a donc  $1  + \left(\dfrac14\right)^n \leqslant u_n \leqslant 2 - \dfrac{n}{n+1}$ avec  $\displaystyle\lim_{n \to + \infty}1 + \left(\dfrac14\right)^n = \displaystyle\lim_{n \to + \infty}2 - \dfrac{n}{n + 1} = 1$. 

D'après le théorème des gendarmes, on peut dire que la suite $(u_n)$ est convergente et que $\displaystyle\lim_{n \to + \infty} u_n=1$.\hfill
Réponse \textbf{b.}

\item Soit $f$ la fonction définie sur $]0~;~+\infty[$ par $f(x) = x^2 \ln x$.










Si $F(x) = \dfrac13 x^3 \left(\ln x - \dfrac13 \right)$, alors 
$F'(x) = x^2\left(\ln x - \dfrac13 \right) + \dfrac13 x^3 \times \dfrac{1}{x} = x^2 \ln x - \dfrac13 x^2  + \dfrac13 x^2 = x^2 \ln x = f(x)$. \hfill Réponse \textbf{a.}

\item 








$2 + \dfrac{3\text{e}^{-x} - 5}{\text{e}^{-x} + 1} = \dfrac{2\text{e}^{-x} + 2  + 3\text{e}^{-x} - 5}{\text{e}^{-x} + 1} = \dfrac{5\text{e}^{-x} - 3}{\text{e}^{-x} + 1}$ ; en multipliant chaque terme par $\text{e}^x$, on obtient : $\dfrac{5 - 3\text{e}^x}{1 + \text{e}^x}$. \hfill Réponse \textbf{a.}
\end{enumerate}

\bigskip
